Automated memory management, commonly referred to as garbage collection, is a core feature of the Java programming language and an integral part of both the Java Runtime Environment (JRE) and the Java Virtual Machine (JVM). In Java, programs are compiled into bytecode and executed on the JVM, where objects are dynamically allocated in a region of memory known as the heap. As execution progresses, many of these objects become unnecessary. The garbage collector is responsible for identifying such unreachable objects and reclaiming their memory, ensuring efficient utilization of system resources. Because of its direct involvement in memory handling, the Garbage Collector (GC) plays a significant role in determining the overall performance and efficiency of an application~\cite{carpen2015performance, lengauer2014taming}. In modern software systems that demand continuous availability and responsiveness, understanding the behavior and impact of different garbage collectors has become increasingly important.

A key advantage of garbage collection in Java is its ability to automatically manage memory by removing unused or unreachable objects. This eliminates the need for developers to manually allocate and deallocate memory, as is required in languages such as C and C++. While manual memory management offers fine-grained control and can sometimes yield performance benefits, it also introduces additional complexity and risk of errors. Despite ongoing debates between these approaches, garbage collection has become a standard feature in many modern programming languages. Furthermore, Java provides various tuning mechanisms that allow developers to optimize garbage collection behavior in cases where it negatively impacts application performance.

With the introduction of Java 13, the JVM supports multiple garbage collectors, including Serial GC, Parallel GC, Concurrent Mark Sweep GC, Epsilon GC, Garbage First GC (G1GC), Shenandoah GC~\cite{flood2016shenandoah}, and Z Garbage Collector (ZGC)~\cite{liden2018zgc}. Among these, G1GC serves as the default collector, while Shenandoah GC and ZGC represent more recent advancements designed to address performance challenges, particularly in large-scale and latency-sensitive applications. These newer collectors have gained attention for their ability to reduce pause times and improve overall system responsiveness.

Performance evaluation of garbage collectors is typically based on two primary metrics: latency and throughput. Latency reflects the responsiveness of an application and is largely influenced by pause events during garbage collection. In certain phases, the garbage collector must halt all application threads to perform essential memory management operations. These pauses, known as Stop The World (STW) events, temporarily prevent the application from making progress and thus contribute to latency. Minimizing both the duration and frequency of such pauses is crucial for latency-sensitive applications. Throughput, on the other hand, measures the amount of useful work performed by an application over time. It is often expressed as the proportion of total execution time not spent on garbage collection or in terms of completed operations within a given time frame. While optimizing for throughput may tolerate longer pauses, real-world systems typically require a careful balance between throughput and latency to achieve optimal performance.

In addition to these traditional performance metrics, energy efficiency has become an increasingly relevant concern in modern computing environments. As large-scale systems continue to grow in complexity and resource demands, understanding the energy footprint of different garbage collection strategies is essential for building sustainable and cost-effective systems.

This challenge is particularly evident in big data systems, which are commonly developed using managed languages such as Java, C\#, and Scala. These systems generate vast numbers of objects while processing large volumes of data, placing significant pressure on memory management mechanisms. The continuous allocation and reclamation of objects can lead to substantial garbage collection overhead, sometimes accounting for a large portion of total execution time. In extreme cases, garbage collection may consume nearly half of the runtime, severely impacting system performance. The problem is further amplified in distributed and latency-sensitive cloud environments, where long pause times on a single node can delay the entire system, leading to poor user experience. Consequently, efficient memory management remains a critical challenge in the context of big data applications.

To address these challenges, modern garbage collectors employ various heuristics and adaptive strategies based on observed application behavior. For instance, Garbage First GC prioritizes the collection of younger objects based on the observation that recently allocated objects are more likely to become unreachable. Similarly, Shenandoah GC uses adaptive heuristics that consider factors such as heap occupancy, allocation rates, and past behavior to determine when to initiate collection cycles. These approaches rely on the assumption that application memory usage patterns exhibit some degree of predictability. The effectiveness of a garbage collector often depends on how well it aligns with the characteristics of the application workload, as no single collector is universally optimal across all scenarios.

In this project, we evaluate and compare three modern garbage collectors, namely Garbage First GC (G1GC), Shenandoah GC, and ZGC. The comparison is performed across multiple benchmarks, with particular attention given to workloads derived from big data libraries. We analyze their performance under varying heap sizes and assess their behavior in terms of latency, throughput, and overall efficiency. We also  perform a detailed tail-latency analysis of the garbage collectors. Additionally, we measure the energy footprint of these garbage collectors to provide a more comprehensive understanding of their impact in real-world computing environments.